\documentclass[11pt]{article}

\usepackage[a4paper,margin=30mm]{geometry}
\usepackage[T1]{fontenc}
\usepackage[utf8]{inputenc}
\usepackage{lmodern}
\usepackage{microtype}
\usepackage{amsmath,amssymb,amsthm,mathtools}
\usepackage{booktabs}
\usepackage{xcolor}
\usepackage[hidelinks]{hyperref}
\usepackage[nameinlink,noabbrev]{cleveref}

\hypersetup{
  pdftitle={Two-slack window-covering words and a triangular age automaton},
  pdfauthor={Nicolo Pandolfo},
  pdfsubject={Exact transfer enumeration of OEIS A242322},
  pdfkeywords={combinatorics on words, finite automata, transfer matrices, OEIS A242322}
}

\newtheorem{theorem}{Theorem}[section]
\newtheorem{proposition}[theorem]{Proposition}
\newtheorem{lemma}[theorem]{Lemma}
\newtheorem{corollary}[theorem]{Corollary}
\theoremstyle{remark}
\newtheorem{remark}[theorem]{Remark}

\newcommand{\T}{\mathcal T}
\newcommand{\E}{\mathcal E}
\newcommand{\one}{\boldsymbol 1}

\title{\textbf{Two-slack window-covering words}\\
\large A triangular last-occurrence automaton for OEIS A242322}
\author{Nicol\`o Pandolfo}
\date{July 23, 2026\\\small Draft for independent mathematical review}

\begin{document}
\maketitle

\begin{abstract}
For integers $n,k\geq1$, OEIS A242322 counts canonical words of length
$n+k+2$ over $\{0,1,\ldots,k\}$ in which every consecutive block of length
$k+3$ contains every symbol.  We encode a valid prefix by the two missing
values among the possible ages of the most recent symbol occurrences.  This
gives an exact transfer matrix $A_k$ with
\[
 \binom{k+2}{2}
\]
states, explaining the triangular column orders observed in the OEIS entry.
If $w_k$ is an explicit initial vector, then
\[
 T(n,k)=w_k^{\mathsf T}A_k^{\,n-1}\one,\qquad
 \sum_{n\geq1}T(n,k)x^n
 =xw_k^{\mathsf T}(I-xA_k)^{-1}\one.
\]
The matrix is unimodular, so its characteristic recurrence has order exactly
$\binom{k+2}{2}$.  We also prove the stable-row formula
\[
 T(n,k)=\frac{a_n}{8}k^4+\frac{b_n}{12}k^3
       +\frac{c_n}{8}k^2+\frac{12n+19}{12}k+1
 \qquad (k\geq\max\{1,n-4\}),
\]
where
\[
\begin{aligned}
 a_n&=4\cdot3^{n-1}-2^{n+1}+1,\\
 b_n&=-4\cdot3^n+(6n+9)2^n-6n-1,\\
 c_n&=4\cdot3^{n-1}+(12-4n)2^n+4n^2-5.
\end{aligned}
\]
Thus the previously empirical triangular recurrences and quartic
stabilization have a common finite-state explanation.
\end{abstract}

\section{Introduction}

Let $\Sigma_k=\{0,1,\ldots,k\}$.  A word over $\Sigma_k$ is
\emph{canonical} if its new symbols are introduced in the order
$0,1,\ldots,k$; equivalently, it is a restricted-growth word.  Define
\begin{equation}\label{eq:def}
 T(n,k)=\#\left\{
 w\in\Sigma_k^{\,n+k+2}:
 \begin{array}{l}
 w\text{ is canonical, and}\\
 \text{every factor of length }k+3\text{ contains }\Sigma_k
 \end{array}\right\}.
\end{equation}
This is the table A242322 \cite{oeisA242322}.  Its entry records empirical
column recurrence orders
\[
 3,6,10,15,21,28
\]
for $k=1,\ldots,6$, and empirical quartic formulas for its first fifteen
rows.  The orders are the triangular numbers $\binom{k+2}{2}$.

The main point of this paper is that the triangular numbers are not a numerical
coincidence.  At any time, record the age of the most recent occurrence of
each symbol.  A valid prefix has $k+1$ distinct ages among the $k+3$ possible
ages $0,1,\ldots,k+2$, and age zero is always present.  Exactly two positive
ages are therefore missing.  Those two missing ages form a complete state
description.

This produces a transfer system substantially smaller than the literal suffix
automaton.  It also retains enough structure to evaluate all fixed rows in the
stable range.  The resulting proofs use only finite-state enumeration, exact
linear algebra, and finite sums.

\section{Relabeling and last-occurrence ages}

Put
\[
 q=k+1,\qquad r=k+2=q+1.
\]
The alphabet has $q=r-1$ symbols and the required window length is $r+1$.

\begin{lemma}[Relabeling orbits]\label{lem:orbits}
Every valid labelled word uses all $q$ symbols.  Relabeling by the symmetric
group $S_q$ acts freely, and every orbit contains exactly one canonical word.
Consequently, canonical words may be counted as labelled words divided by
$q!$.
\end{lemma}

\begin{proof}
The first window contains every symbol.  A relabeling that fixes the word
therefore fixes all symbols, so the action is free.  Relabeling the symbols by
their order of first appearance produces the unique canonical representative
of each orbit.
\end{proof}

For a valid prefix ending at position $t$, let
\[
 \alpha_t(s)=t-\max\{u\leq t:w_u=s\}
\]
be the \emph{age} of symbol $s$.  The ages are distinct, age zero is present,
and validity is equivalent to $\alpha_t(s)\leq r$ for every symbol.

\begin{lemma}[Triangular state space]\label{lem:states}
After the first complete window, the age set has the form
\[
 \{0,1,\ldots,r\}\setminus\{i,j\},
 \qquad 1\leq i<j\leq r.
\]
Thus the quotient state space is
\[
 \E_r=\{(i,j):1\leq i<j\leq r\},
 \qquad |\E_r|=\binom r2=\binom{k+2}{2}.
\]
\end{lemma}

\begin{proof}
There are $q=r-1$ symbols, hence $r-1$ distinct ages.  They lie among the
$r+1$ integers $0,\ldots,r$, and zero occurs as the age of the final symbol.
Their complement therefore consists of exactly two positive integers.
\end{proof}

\begin{lemma}[Transitions]\label{lem:transitions}
The successors of $(i,j)\in\E_r$ are as follows.
\begin{enumerate}
\item If $j<r$, the unique successor is $(i+1,j+1)$.
\item If $j=r$, put $v=i+1$.  The successors are all pairs incident with $v$:
\[
 \bigl\{(\min\{v,u\},\max\{v,u\}):1\leq u\leq r,\ u\neq v\bigr\}.
\]
\end{enumerate}
Every listed edge is realized by exactly one appended symbol.
\end{lemma}

\begin{proof}
If $j<r$, age $r$ is present.  Its symbol must be appended, since otherwise
its age would become $r+1$.  Resetting that age to zero and increasing all
other ages shifts the two missing ages to $i+1,j+1$.

If $j=r$, all current ages are at most $r-1$, so any symbol may be appended.
Its current age $a$ ranges over $\{0,\ldots,r-1\}\setminus\{i\}$.  After the
update, the missing ages are $i+1$ and $a+1$.  As $a$ varies, $a+1$ takes
every value in $\{1,\ldots,r\}\setminus\{i+1\}$ exactly once.
\end{proof}

Let $A_k$ be the $\binom r2$ square matrix indexed by $\E_r$, with
\[
 (A_k)_{e,f}=
 \begin{cases}
 1,&e\longrightarrow f\text{ in \cref{lem:transitions}},\\
 0,&\text{otherwise}.
 \end{cases}
\]

\begin{lemma}[Initial vector]\label{lem:initial}
The number of canonical first windows whose terminal age state is $(i,j)$ is
\[
 w_k(i,j)=i(j-1).
\]
\end{lemma}

\begin{proof}
A first window has length $r+1$ and contains all $q=r-1$ symbols.  Fix the
set of last-occurrence ages, namely $\{0,\ldots,r\}\setminus\{i,j\}$.
Assigning the $q$ labels to these last-occurrence positions contributes $q!$.
At the non-last position of age $i$, its symbol may have any of the $i$ last
occurrences of smaller age.  At age $j$, there are $j-1$ choices, since age
$i$ is not a last-occurrence position.  These choices are independent.
There are therefore $q!i(j-1)$ labelled windows.  Divide by $q!$ using
\cref{lem:orbits}.
\end{proof}

\section{Exact transfer formula and column recurrences}

\begin{theorem}[Exact transfer enumeration]\label{thm:transfer}
For all $n,k\geq1$,
\begin{equation}\label{eq:transfer}
 T(n,k)=w_k^{\mathsf T}A_k^{\,n-1}\one,
\end{equation}
where $w_k(i,j)=i(j-1)$ and $\one$ is the all-one column vector.  Hence
\begin{equation}\label{eq:gf}
 \mathcal T_k(x):=\sum_{n\geq1}T(n,k)x^n
 =xw_k^{\mathsf T}(I-xA_k)^{-1}\one.
\end{equation}
\end{theorem}

\begin{proof}
\Cref{lem:initial} counts the possible first windows by state.
\Cref{lem:transitions} gives a one-to-one correspondence between appended
symbols and state edges.  A word counted by $T(n,k)$ has $n-1$ transitions
after its first window, and its final state is unrestricted.  This proves
\eqref{eq:transfer}; summing the matrix geometric series proves \eqref{eq:gf}.
\end{proof}

Define
\[
 D_k(x)=\det(I-xA_k)=\sum_{\ell=0}^{N}d_{k,\ell}x^\ell,
 \qquad N=\binom{k+2}{2}.
\]

\begin{proposition}[Unimodularity]\label{prop:unimodular}
The matrix $A_k$ is unimodular.  In particular,
$d_{k,0}=1$, $|d_{k,N}|=1$, and $D_k$ has degree exactly $N$.
\end{proposition}

\begin{proof}
Order first the rows $(i,j)$ with $j<r$, then the boundary rows $(i,r)$.
Order first the columns $(u,v)$ with $2\leq u<v\leq r$, then the columns
$(1,v)$.

The interior transition $(i,j)\mapsto(i+1,j+1)$ is a bijection from the
interior rows to the first group of columns.  The corresponding block is a
permutation matrix, and the upper-right block is zero.  In boundary row
$(i,r)$, the only target in the second column group is $(1,i+1)$.  The
lower-right block is therefore an identity matrix.  The reordered matrix is
block lower triangular with unimodular diagonal blocks.
\end{proof}

\begin{corollary}[Triangular-order recurrence]\label{cor:recurrence}
For every $k\geq1$, the column $T(n,k)$ satisfies
\begin{equation}\label{eq:recurrence}
 \sum_{\ell=0}^{N}d_{k,\ell}T(n-\ell,k)=0
 \qquad(n\geq N+1),
 \qquad N=\binom{k+2}{2}.
\end{equation}
Thus A242322 has a characteristic recurrence of the triangular order displayed
empirically on OEIS.
\end{corollary}

\begin{proof}
Cayley--Hamilton gives
$\sum_{\ell=0}^{N}d_{k,\ell}A_k^{N-\ell}=0$.
Multiplication by the initial and terminal vectors gives
\eqref{eq:recurrence}.  The final coefficient is nonzero by
\cref{prop:unimodular}.
\end{proof}

\begin{remark}
The result asserts a recurrence with nonzero endpoints and order $N$; it does
not require a claim that $N$ is the minimal scalar order for every $k$.
Exact computation shows no scalar cancellation for $1\leq k\leq10$.
\end{remark}

For example, exact characteristic-polynomial calculation gives
\begin{align*}
 D_1(x)&=1-x-x^2-x^3,\\
 D_2(x)&=1-x-2x^2-2x^3-2x^4+x^5+x^6,\\
 D_3(x)&=1-x-2x^2-3x^3-5x^4-6x^5+x^6+x^7+x^9+x^{10}.
\end{align*}
The first two agree with the recurrences printed on A242322.

\section{Stable rows}

We now evaluate \eqref{eq:transfer} when the alphabet is sufficiently large
relative to the number of transitions.

Use distance coordinates
\[
 a=r-j,\qquad b=r-i,\qquad 0\leq a<b\leq r-1.
\]
Let $H_m(a,b)$ be the number of length-$m$ walks starting at this state, with
unrestricted endpoint, and put $G_m(b)=H_m(0,b)$.  Then $H_0(a,b)=1$, and the
transition rules give
\begin{align}
 H_m(a,b)&=H_{m-1}(a-1,b-1) &&(a>0),\label{eq:Hinterior}\\
 G_m(b)&=\sum_{\substack{0\leq c\leq r-1\\c\neq b-1}}
 H_{m-1}\bigl(\min\{b-1,c\},\max\{b-1,c\}\bigr).
 \label{eq:Gboundary}
\end{align}

\begin{lemma}[Boundary walks]\label{lem:boundary}
Suppose $s\geq1$ and $r\geq s-1$.  For every $1\leq b\leq r-1$,
\[
 G_s(b)=
 \begin{cases}
 2^{s-1}(r-2)+1,&b\geq s,\\[2mm]
 2^{b-1}3^{p-1}r^2+
 \bigl(2^{p-1}+2^{s-1}-5\cdot2^{b-1}3^{p-1}\bigr)r\\
 \qquad{}+2^b3^p-2^s-2^p+1,
 &b<s,\quad p=s-b.
 \end{cases}
\]
\end{lemma}

\begin{proof}
We use strong induction on $s$.  Repeated application of
\eqref{eq:Hinterior} gives
\[
 H_m(a,b)=
 \begin{cases}
 1,&m\leq a,\\
 G_{m-a}(b-a),&m>a.
 \end{cases}
\]
Substitution into \eqref{eq:Gboundary} yields, for $b<s$ and $p=s-b$,
\begin{equation}\label{eq:Grec}
 G_s(b)=\sum_{e=1}^{b-1}G_{p+e}(e)
       +\sum_{e=1}^{r-b}G_p(e).
\end{equation}
For $b\geq s$, the same decomposition has only the constant branch and gives
\[
 G_s(b)=r-1+\sum_{t=1}^{s-1}
 \bigl(2^{t-1}(r-2)+1-1\bigr)
 =2^{s-1}(r-2)+1.
\]

For the second sum in \eqref{eq:Grec}, the induction hypothesis and geometric
sums give, for every $R\geq p-1$,
\[
 \sum_{e=1}^{R}G_p(e)
 =R\bigl(2^{p-1}(r-2)+1\bigr)+U_p(r),
\]
where
\[
\begin{aligned}
 U_p(r)={}&(3^{p-1}-2^{p-1})r^2\\
 &+(-5\cdot3^{p-1}+3\cdot2^p-1)r
   +2\cdot3^p-4\cdot2^p+2.
\end{aligned}
\]
The first sum in \eqref{eq:Grec} uses the nonconstant induction branch.
Summing its three coefficients over $e=1,\ldots,b-1$, then adding the preceding
display with $R=r-b$, gives exactly the stated expression.  The same
calculation remains valid at the limiting value $r=s-1$, where $R=p-1$.
\end{proof}

\begin{theorem}[Stable quartic rows]\label{thm:rows}
For $n,k\geq1$ with $k\geq\max\{1,n-4\}$,
\begin{equation}\label{eq:row}
 T(n,k)=\frac{a_n}{8}k^4+\frac{b_n}{12}k^3
       +\frac{c_n}{8}k^2+\frac{12n+19}{12}k+1,
\end{equation}
where
\begin{align*}
 a_n&=4\cdot3^{n-1}-2^{n+1}+1,\\
 b_n&=-4\cdot3^n+(6n+9)2^n-6n-1,\\
 c_n&=4\cdot3^{n-1}+(12-4n)2^n+4n^2-5.
\end{align*}
The uniform boundary cannot be lowered: at $(n,k)=(6,1)$ the true value is
$149$, whereas the polynomial gives $151$.
\end{theorem}

\begin{proof}
Put $m=n-1$.  In distance coordinates, the initial weight from
\cref{lem:initial} is
\[
 (r-b)(r-a-1).
\]
Consequently,
\begin{equation}\label{eq:rowsum}
 T(n,k)=\sum_{0\leq a<b\leq r-1}
 (r-b)(r-a-1)H_m(a,b).
\end{equation}
The hypothesis $k\geq n-4$ is exactly $r\geq m-1$, so
\cref{lem:boundary} applies to every nontrivial term after
\eqref{eq:Hinterior} is iterated.  Substitute that lemma into
\eqref{eq:rowsum}, split the sum at $b-a=m-a$, and use
\[
 \sum_{u=1}^{M}u=\frac{M(M+1)}2,\quad
 \sum_{u=1}^{M}u^2=\frac{M(M+1)(2M+1)}6,\quad
 \sum_{u=1}^{M}2^u=2^{M+1}-2.
\]
Collecting powers of $k=r-2$ gives, respectively,
\[
 \frac{a_n}{8},\quad\frac{b_n}{12},\quad
 \frac{c_n}{8},\quad\frac{12n+19}{12},\quad1,
\]
which proves \eqref{eq:row}.  Direct enumeration at $(6,1)$ gives the final
boundary statement.
\end{proof}

\section{Verification, scope, and disclosure}

No proof step depends on fitting numerical data.  Two standard-library-only
programs accompany the manuscript.  The feasibility verifier implements an
equality-pattern quotient, a raw labelled suffix DFA, and literal
restricted-growth enumeration.  The theorem verifier independently implements
the age automaton, exact characteristic polynomials, the boundary-walk lemma,
and the stable-row formula.  Together they check:
\begin{enumerate}
\item the triangular state classification and every transition;
\item agreement of four enumeration routes;
\item the first $10\times10$ values displayed on A242322;
\item characteristic recurrences and unimodularity through $k=10$;
\item the boundary-walk formula, including its limiting case;
\item the stable-row theorem by matrix walks through $n=40$, and independently
      by direct summation of the boundary lemma through $n=100$.
\end{enumerate}
Finite checks support the audit but do not replace the proofs.

The matrix formula is an exact generating function for every column.  The
paper does not claim a closed coefficient formula for $D_k(x)$, nor does it
claim that its degree is always the minimal scalar recurrence order.  The
literature search found the empirical statements on A242322 but no prior proof
of the age-state reduction; this is a qualified status report rather than an
absolute priority claim.

AI systems were used during exploration, adversarial checking, and drafting.
The mathematical evidence consists of the explicit arguments and deterministic
verification programs, not AI output.  This draft is marked for independent
mathematical review.

\section*{Acknowledgment}
The definition, numerical table, and empirical formulas originate with R. H.
Hardin and the contributors to OEIS A242322.

\begin{thebibliography}{9}

\bibitem{oeisA242322}
OEIS Foundation Inc.,
\emph{A242322: Window-covering canonical arrays with two extra positions},
\url{https://oeis.org/A242322}, accessed July 23, 2026.

\bibitem{stanley}
R. P. Stanley,
\emph{Enumerative Combinatorics, Volume 1}, 2nd ed.,
Cambridge University Press, 2012, Section 4.7.

\bibitem{flajolet}
P. Flajolet and R. Sedgewick,
\emph{Analytic Combinatorics},
Cambridge University Press, 2009, Chapters I and V.

\end{thebibliography}

\end{document}
